\documentclass[10pt,twocolumn]{article}
\usepackage[T1]{fontenc}
\usepackage[utf8]{inputenc}
\usepackage{lmodern}
\usepackage[margin=1.8cm,top=2.4cm,bottom=2cm,columnsep=0.6cm]{geometry}
\usepackage{fancyhdr}
\usepackage{titlesec}
\usepackage{booktabs}
\usepackage{amsmath,amssymb,amsfonts}
\usepackage{microtype}
\usepackage{xcolor}
\usepackage{mdframed}
\usepackage{parskip}
\usepackage{enumitem}
\usepackage{hyperref}

\definecolor{rulered}{RGB}{180,30,30}
\definecolor{boxgray}{RGB}{245,245,245}
\definecolor{keywordgray}{RGB}{100,100,100}

\pagestyle{fancy}
\fancyhf{}
\fancyhead[L]{\textsc{\small Claude Mag}}
\fancyhead[C]{\small\textit{Machine Learning Research Digest}}
\fancyhead[R]{\small 2026-10-19}
\fancyfoot[C]{\small\thepage}
\renewcommand{\headrulewidth}{0.8pt}

\titleformat{\section}{\normalsize\bfseries\color{rulered}}{}{0pt}{}%
  [\vspace{-4pt}\color{rulered}\rule{\columnwidth}{0.4pt}]
\titlespacing{\section}{0pt}{8pt}{4pt}

\hypersetup{colorlinks=true,urlcolor=rulered,linkcolor=black}

\begin{document}
\setlength{\parindent}{0pt}
\setlength{\parskip}{4pt}

{\small\color{keywordgray}\textbf{Keywords:}
  agentic video \textbullet{}
  world-state memory \textbullet{}
  cross-shot consistency \textbullet{}
  video generation}\par\vspace{4pt}

{\LARGE\bfseries\raggedright
  Any Time, Any Space}\par\vspace{4pt}

{\large\itshape\raggedright
  A World-State-Centric Framework for Agentic Video Generation Maintains
  Two Persistent Memories---Scene Geometry and Indexed Dynamic Events---to
  Ensure Full Spatio-Temporal Coherence Across Nonconsecutive Shots}\par\vspace{6pt}

{\small\textbf{By} Weijie Wang, Haomin Zhang, Bohan Zhuang et al.}\par\vspace{2pt}

{\small\color{keywordgray}arXiv:2610.20571 \textbullet{} October 2026}
\par\vspace{2pt}\noindent{\color{rulered}\rule{\columnwidth}{0.6pt}}\vspace{6pt}

Long-form narrative video generation requires that any shot---requested at any
time, depicting any space---remain spatially and temporally consistent with every
other shot in the production. ATAS (\emph{Any Time, Any Space}) achieves this
by maintaining a \emph{scene memory} of registered static geometry and an
\emph{event memory} of dynamic observations indexed by scene, event, and time.
Before generating each shot the system retrieves the relevant evidence; after
generation a review gate accepts or rejects the clip before registering it into
memory. Across four video generators ATAS improves every reported cross-shot
consistency metric and tops all six human-rating criteria.

\begin{mdframed}[backgroundcolor=boxgray,linecolor=rulered,linewidth=1pt,
    innerleftmargin=6pt,innerrightmargin=6pt,
    innertopmargin=6pt,innerbottommargin=6pt]
\textbf{\small AT A GLANCE}\par\vspace{4pt}
\begin{description}[leftmargin=0pt,labelsep=4pt,itemsep=2pt,parsep=0pt,
    font=\small\bfseries]
  \item[Problem:] Long-form agentic video suffers from location drift and temporal
    contradictions when shots are generated nonconsecutively; existing methods
    lack a shared spatial and temporal state.
  \item[Why it matters:] Agentic video is an emerging AI application where
    coherence failures break narrative trust; no prior approach solves
    multi-shot spatio-temporal consistency.
  \item[Method:] Scene memory stores registered geometry per location; event
    memory stores dynamic observations indexed by scene, event, and time;
    a review gate decouples film delivery from memory admission.
  \item[Result:] Highest aggregate benchmark score per backbone; tops all six
    human-rating criteria across four video generators.
  \item[Evidence:] Four video generators; spatial revisit, event recapture, and
    temporal consistency benchmarks; human evaluation on six criteria.
\end{description}
\end{mdframed}

\section{The Big Picture}

Narrative video pipelines built from large generative models suffer from
\emph{scene drift}: ask the same model to re-enter a room from a new camera
angle three shots later, and it will produce a room that contradicts the
layout established earlier. The same problem afflicts event recapture---a
flashback to an event the viewer already witnessed will not match the original.

These failures arise because every shot is generated independently. The model
conditions on a text prompt and, optionally, the most recent frame, but neither
source encodes the durable spatial and temporal facts the narrative has
established. ATAS resolves this by replacing the stream of frames with a
\emph{world-state representation}: two structured memories that persist across
the entire production, regardless of generation order.

\section{Why Existing Approaches Fall Short}

\textbf{Prompt conditioning} passes shot descriptions between generation calls.
It provides no spatial grounding; the model must re-invent geometry from text at
each step, and two independent descriptions of the same room will produce
inconsistent geometry.

\textbf{Recent-frame conditioning} threads the previous clip's last frame into
the current generation. This provides visual continuity for adjacent shots, but
cannot support true out-of-order revisits or event recall across large temporal
distances in the screenplay.

\textbf{3D-guided methods} impose rigid scene constraints extracted from
monocular depth or NeRF reconstruction. These are difficult to update as the
narrative introduces new objects, character motion, and destructive events; they
also require per-production reconstruction passes that scale poorly.

\textbf{Agent planning without state} improves shot ordering and prompt
engineering but leaves individual generation calls unmodified; without a
persistent world model, cross-shot coherence is not guaranteed.

\section{Core Mechanism: ATAS}

ATAS structures the agentic loop around two memories:

\textbf{Scene Memory} stores a geometric descriptor $G(s_i)$ for each named
scene $s_i$. When a location is first visited, the generated clip is
reconstructed and the geometry registered. All subsequent shots at $s_i$
retrieve $G(s_i)$ and condition on it, ensuring consistent spatial layout
regardless of camera angle or narrative time.

\textbf{Event Memory} stores a clip descriptor $D(s_i, e_j, t)$ for each
combination of scene $s_i$, event $e_j$, and time $t$. When an action or
state change is generated, the resulting clip is reconstructed and registered
at the appropriate indices. A later shot requesting a replay retrieves the
stored entry rather than re-generating from a text description.

\textbf{Review Gate}: Before registering any clip, ATAS applies quality
and state-consistency checks. A rejected clip is regenerated; film delivery is
independent of memory admission, so the output stream never blocks and
hallucinated geometry is never registered.

\section{Technical Formulation}

Let $\mathcal{S} = \{s_1, \ldots, s_N\}$ be named scenes and
$\mathcal{E} = \{e_1, \ldots, e_M\}$ be narrative events. Scene memory maps
each scene to its registered geometry:
\[
  G(s_i) = \mathrm{Reconstruct}\!\left(v^{(1)}(s_i)\right)
\]
where $v^{(1)}(s_i)$ is the first accepted clip at scene $s_i$.

Event memory maps a triple to a dynamic descriptor:
\[
  D: \mathcal{S} \times \mathcal{E} \times \mathbb{R}_{\geq 0} \to \mathcal{D}
\]

For a shot request $r = (s_i, e_j, t, \mathrm{cam})$, the conditioning evidence
is the concatenation of the two retrieved descriptors:
\[
  C(r) = G(s_i) \oplus D(s_i, e_j, t)
\]
and the generator produces:
\[
  v_r = \mathrm{VideoGen}(r,\; C(r))
\]

After the review gate accepts $v_r$, the event memory is updated:
\[
  D(s_i, e_j, t) \;\leftarrow\; \mathrm{Reconstruct}(v_r)
\]

If the review gate rejects $v_r$, the clip is regenerated; the memory state
is unchanged, preventing error propagation.

\section{Learning or Inference Procedure}

\begin{enumerate}[itemsep=2pt,parsep=0pt]
  \item Parse the screenplay into shot requests $r = (s_i, e_j, t, \mathrm{cam})$.
  \item For each request, query scene memory for $G(s_i)$ and event memory for
    $D(s_i, e_j, t)$; retrieve empty descriptors for unseen entries.
  \item Condition the video generator on $C(r)$ and generate clip $v_r$.
  \item Apply the review gate; if accepted, update $D(s_i, e_j, t)$ and
    $G(s_i)$ (on first visit); deliver $v_r$ to the output stream.
  \item If rejected, return to step 3 with a refreshed generation seed.
  \item Repeat until all shots are delivered.
\end{enumerate}

No additional model training is required. ATAS is a retrieval-and-conditioning
framework that wraps any video generation backbone.

\section{What the Guarantee Says}

ATAS does not make formal guarantees on per-shot quality, which depends on the
underlying generator. Its empirical claims are:

\begin{itemize}[itemsep=2pt,parsep=0pt]
  \item Every reported cross-shot consistency metric improves across all four
    tested video generators when ATAS is added.
  \item ATAS achieves the highest aggregate benchmark score per backbone on the
    spatial-revisit and event-recapture tasks.
  \item Human raters score ATAS highest on all six criteria (narrative coherence,
    spatial accuracy, temporal consistency, visual quality, event fidelity,
    audio-visual alignment) across all tested generators.
\end{itemize}

\section{Experimental Findings}

\begin{itemize}[itemsep=2pt,parsep=0pt]
  \item \textbf{Backbones:} Four video generators covering diffusion-based and
    autoregressive architectures.
  \item \textbf{Spatial revisit:} ATAS outperforms compared agent-based and
    3D-guided methods on the target-image spatial revisit task.
  \item \textbf{Event recapture:} ATAS leads on three event-recapture criteria;
    prior methods degrade to text-driven re-generation.
  \item \textbf{Temporal consistency:} All temporal-consistency metrics improve
    under ATAS for every backbone.
  \item \textbf{Human evaluation:} ATAS tops all six human-rating criteria
    across all tested generators; previous methods trail on at least two.
  \item \textbf{Out-of-order generation:} ATAS maintains consistency even when
    shots are generated out of screenplay order.
\end{itemize}

\section{Ablations and Interpretation}

\textbf{Without scene memory:} Spatial layout drifts across revisit shots;
objects disappear or change position; the spatial-revisit score drops sharply.

\textbf{Without event memory:} Event recapture degrades to text-driven
re-generation; temporal positions shift; event-fidelity scores drop.

\textbf{Without review gate:} Hallucinated geometry is registered and
propagates to subsequent shots; errors compound as the production grows.

\textbf{Memory granularity:} Finer time indexing in event memory yields better
event recapture; coarser indexing reduces memory overhead at the cost of
temporal precision.

\textbf{Why retrieval beats prompting:} Geometry extracted from accepted clips
encodes exact object positions that a text description cannot; conditioning on
this evidence directly constrains the generator's spatial output distribution,
while text descriptions leave spatial degrees of freedom underdetermined.

\section*{Reference}

Weijie Wang, Haomin Zhang, Bohan Zhuang, et al.
\textbf{Any Time, Any Space: A World-State-Centric Framework for Agentic Video
Generation}.
arXiv:2610.20571, October 2026.\\
\url{https://arxiv.org/abs/2610.20571}

\end{document}
